\section{Introduction}\label{sec:current-overview}\label{sec:introduction}

A predictive state records the information in an experimental past that
can affect future laws.  Its linear dimension measures one form of
compression.  Implementing those laws with finite positive memory also
requires stochastic transitions, a common decoder and an admissible
initialization.  Implementing an optimal policy introduces a further
question: how much does the best attainable cost change when the device
is allowed one additional persistent label?  Predictive dimension,
the number of physical states and the number of optimal decision values
need not answer that question.

This paper relates predictive, convex and spectral structure to finite
positive implementations of prepared causal experiments.  A controller
has one fixed parameter-independent program.  It sees only the declared
reports, and retains no uncounted history, clock or private seed.
Persistent labels, the protected physical state, program-table bits and
working arithmetic storage are separate resources.  A prescribed pure
initial label is counted.  Freely prepared mixed starts are identified
whenever that different convention is used.

The central result is an exact finite-memory variational principle for
reset instruments.  It covers every memory budget, singular failure
effects, arbitrary finite Hilbert dimension, several commands,
state-dependent costs and failure probabilities, and tied optimal
actions.  The proof removes arbitrary stochastic memory updates without
adding a row.  An explicit family then solves every memory budget and
has arbitrarily long strict resource hierarchies.  These results are
developed alongside the general shared-preparation realization theory,
operational approximation lower bounds and the original causal
comparison.  The theorem classes and their quantifiers remain distinct.

\subsection*{The complete reset memory frontier}

Let $a\in\mathcal A$ be a command and $i\in\{1,\ldots,s\}$ a failure
report.  The physical failure branch is
\[
 \mathcal I_{a,i}(\rho)=\tr(E_{a,i}\rho)\sigma_i,
 \qquad E_{a,i}\succeq0,\quad \sum_iE_{a,i}\preccurlyeq I.
\]
The report itself prepares the known state $\sigma_i$.  The observer
does not receive an additional state measurement or free reset.  A
successful call must terminate and failure forbids termination.  The
mean call cost is $\tr(G_a\rho)$ with
$c_*I\preccurlyeq G_a\preccurlyeq c^*I$ and $c_*>0$; a realized cost is
not an extra observation.  The initial density matrix is arbitrary and
unknown.  Write $V_M$ for the least worst-initial-state expected cost
over autonomous controllers with at most $M$ live labels and a prescribed
pure initial label.

Absence of a closed failure set is necessary and sufficient for a
finite-cost controller, even with one randomized command row
(Theorem~\ref{thm:reset-control}).  Put
\[
 Q_{iaj}=\tr(E_{a,j}\sigma_i),\qquad c_{ia}=\tr(G_a\sigma_i).
\]
For a row palette $p^1,\ldots,p^M\in\Delta(\mathcal A)$ arising from
a controller with finite worst-initial-state cost, consider the finite
decision problem whose actions are exactly these rows.  Its
least Bellman potential is
\[
 x_i=\min_{1\le k\le M}
           \sum_a p_a^k\left(c_{ia}+\sum_jQ_{iaj}x_j\right).
\]
Theorem~\ref{thm:r78-palette-normal-form} proves that every finite-cost
controller is dominated, separately at every initial density matrix,
by one with the same rows, the same initial row $p^1$, and deterministic
report-only routing $i\mapsto\kappa(i)$.  The argument uses the
all-initial-state requirement at an essential point: starting the
original controller with its fixed initial label from each $\sigma_i$
gives a finite comparison policy for every state of the row-palette
problem.  It does not assert transience of unreachable pairs in the
original product of physical and controller states.

Let $C_0$ be an explicit finite worst-state cost of a full-support
one-row controller.  Theorem~\ref{thm:r78-reset-frontier} gives
\begin{equation}\label{eq:r78-intro-frontier}
\begin{split}
 V_M=\min\{b:\;&p^k\in\Delta(\mathcal A),\quad
      \kappa:\{1,\ldots,s\}\to\{1,\ldots,M\},\quad
      x\in[0,C_0]^s,\quad b\in[0,C_0],\\
 &x_i\ge\sum_a p_a^{\kappa(i)}
                  (c_{ia}+\sum_jQ_{iaj}x_j)\quad(1\le i\le s),\\
 &\sum_a p_a^1(G_a+\sum_i x_iE_{a,i})\preccurlyeq bI\}.
\end{split}
\end{equation}
The minimum is attained and each feasible point constructs an actual
controller.  Conversely, every cost bound below $C_0$ realized by any
stochastic controller gives such a point.  Only $s$ scalar continuation
potentials are needed, independently of $M$; no $sM$-state product
resolvent is imposed.  Singular effects and arbitrary Bellman ties are
allowed.  For exact algebraic primitive data, this compact finite union
of semialgebraic sets decides every cost threshold and gives an exact
minimizer.  No polynomial running time or executed general elimination
is asserted.

Proposition~\ref{prop:r78-frontier-stability} proves continuity of each
$V_M$ in the admissible primitive instrument.  The lower argument uses
the bounded potentials and compact row palettes; the upper argument
uses stability of a fixed proper realizing controller.  Thus a strict
finite resource gap persists under sufficiently small perturbations,
including perturbations that remove commutation among effects and
cost operators.

There are exact growing-dimensional instances.  For any nonempty proper
action subsets $A_1,\ldots,A_s\subset\{1,\ldots,q\}$, the inherited
rational construction has $E_{a,i}=(r/s)I$ and costs in $[1,4]$.
Theorem~\ref{thm:r78-cover-frontier} computes $V_M$ by a finite minimum
over at most $M-1$ continuing actions and a nonempty initial support.
Its optimum has pure continuing rows and an initial row uniform on
that support.  This purity conclusion is for this displayed family;
the general normal form preserves arbitrary mixed command rows.
For $A_i=\{i\}$, $s=q=n\ge2$ and $r=1/2$, the full formula becomes
\begin{equation}\label{eq:r78-intro-hierarchy}
 V_M=4+\frac1n+
      \frac{(n-1)\max\{n-M+1,0\}}{n^2},\qquad M\ge1.
\end{equation}
It gives $V_1>V_2>\cdots>V_{n+1}=V_*$ in dimension $2n$, with
one exact decrement $(n-1)/n^2$ at every step.  In this family the full
gap is exactly twice the older first-residual certificate $D_M$.
The distinction between a lower certificate and the true constrained
optimum is therefore quantitative, not only terminological.

\subsection*{The optimum threshold and positive exposure}

The unrestricted reset optimum is the spectral minimum
$V_*=\min_p\lambda_{\max}(\sum_a p_aL_a)$, where
$L_a=G_a+\sum_i v_iE_{a,i}$ and $v$ is the primitive Bellman optimum.
When one attaining dual preparation $\rho_*$ satisfies
$\tr(\rho_*E_{a,i})>0$ for every command $a$ and report $i$,
the optimal-action sets $\mathcal A_i^*$ give a simpler exact memory
threshold.  Write $\mathcal P_*$ for the optimal initial rows,
$I(p)=\{i:\supp p\subseteq\mathcal A_i^*\}$, and $t(J)$ for the
least number of actions meeting all $\mathcal A_i^*$ with $i\in J$.
The inherited Theorem~\ref{thm:r77-reset-cover} states
\begin{equation}\label{eq:r77-intro-reset}
 m_*=1+\min_{p\in\mathcal P_*}
                 t(\{1,\ldots,s\}\setminus I(p)).
\end{equation}
Positive definite effects are a direct sufficient exposure hypothesis.
The explicit semidefinite example has a one-label optimum that this
cover formula would overcount without exposure.  It remains a valid
counterexample to removing that hypothesis; it is also correctly
handled by the new general frontier.  The arbitrary-cover realization
and its NP-hardness consequence are inherited unchanged.  Their
unrestricted value may be simple even though memory synthesis is hard.

\subsection*{Shared preparations and causal response realization}

Fix a finite group $G$ and finite response laws $b_i:G\to\Delta(C)$.
All preparations must use one command program and one decoder; the
decoder does not receive the preparation identity for free.  For
closed convex constraints on every command word,
Theorem~\ref{thm:r77-shared-normal-form} gives a common permutation
core, including exact charges for prescribed pure initializers.
For the distinct target tables $\mathcal B$, let $P$ range over
$G$-invariant polytopes in the full response simplex containing
$\mathcal B$, and let $v(P)$ count their vertices.  Then
Theorem~\ref{thm:r77-shared-polytope} gives
\begin{equation}\label{eq:r77-intro-polytope}
 \kappa_G(\mathcal B)=\min_P v(P),\qquad
 m_G(\mathcal B)=\min_P\{v(P)+|\mathcal B\setminus\operatorname{vert}P|\}.
\end{equation}
The polytope may extend beyond the target span.  These are different
resource conventions: a freely prepared mixed row and a prescribed
pure device label have different initialization costs.

Horizontal concatenation of the preparation Fourier blocks gives the
joint, potentially noncyclic response module.  With a common group
mean, its permutation multiplicity cover is the exact sufficiently
small positive-contrast memory, with an explicit simultaneous positive
decoder.  Weighted singular values give robust lower certificates.
Exact families on the fixed group $C_2$ require $2q$ shared labels
although every preparation separately needs two, or have mixed memory
two and unbounded pure initialization overhead.  Unequal means have
a nonnegative-rank obstruction.  The full finite-group normal form,
nonabelian multiplicity, quotient duality, and the displayed complete
symmetric-group mean--contrast law retain their proofs and scopes.

\subsection*{Operational approximation and the qubit boundary}

Finite causally prepared heralds followed by delayed terminal tests
give a robust nonnegative-factorization obstruction without requiring
a physical density matrix at each hidden simulator label.  The
decoder may have arbitrary stochastic test columns, but cannot see an
uncounted copy of the earlier history.  Theorem~\ref{thm:r77-operational-antichain}
first controls joint herald and test events, and only then conditions.
For the original qubit phase $z=(3+4i)/5$,
Theorem~\ref{thm:r77-unrestricted-qubit} implies
\begin{equation}\label{eq:r77-intro-arithmetic}
 \delta\le\frac{25\zeta}{128\cdot200^J J^2}
 \quad\Longrightarrow\quad N_{\rm cl}(\delta;B)\ge
 \min\{m:\tbinom m{\lfloor m/2\rfloor}\ge J\},
 \qquad\zeta=\min\{1,7B/8\}.
\end{equation}
At each fixed $B>0$ the unrestricted lower is of order
$\log\log(1/\delta)$; the retained constructive upper is logarithmic.
The matching leading logarithm for physical density dictionaries keeps
its Born-decoder and badly-approximable-phase hypotheses.  These rates
concern stopped response approximation, not the original general-$M$
squared-sign decision-risk frontier.  That separate problem remains
open and is not replaced by the exact reset cost hierarchy.

\subsection*{Classical ingredients, priority and the proof architecture}

Bellman verification and stochastic-shortest-path theory are classical
\cite{BertsekasTsitsiklis91}.  Finite stochastic controller search,
bounded policy iteration and nonlinear optimization of fixed controller
size have an established literature
\cite{Meuleau99,PoupartBoutilier03,AmatoBernsteinZilberstein10}.
Neither writing a nonlinear program nor invoking real algebraic
elimination is claimed as a new discovery.  The specific structural
content here is simultaneous domination with the same counted rows,
removal of all stochastic routing for the all-state reset problem,
the $s$-potential compact certificate in the undiscounted model,
and the full exact resource law and stable hierarchy.

Residual automata, positive realization, representation theory,
equivariant factorizations and antichains retain their classical
attribution \cite{RabinScott59,Rabin63,Heller65,CarlylePaz71,
BenvenutiFarina04,Schwarz64,Flor69,Serre77,GPT13,VandaeleGillisGlineur17}.
Set-cover hardness is used through its explicit physical reduction
\cite{Karp72}.  Targeted primary-source comparisons distinguish these
ingredients from the theorem statements; they do not certify
exhaustive priority or independent editorial significance.

The reset frontier depends on the experiment semantics and the finite
reset existence theorem.  It does not require the group realization or
qubit arithmetic results.  The shared-preparation chain uses the
finite-group convex normal form and its Fourier premises.  The
approximation chain uses operational factorization and phase separation.
The general actual-history comparison and paid construction remain in
the main article; their lower witnesses use the competitor's actual
histories.  The supplement preserves full independent proofs for
finite-order quantum measurement, two-stage reset, memory and
communication, rank-constrained discrimination, performed-law
compactness, variational and dual attainment, spectral memory and local
acquired geometry.  The R58 companion remains independently compilable.
No model-specific value formula is promoted to a solution of arbitrary
causal feedback, global geometry, or the independent A1/A2 programme.
